\documentclass[a4paper]{article}
\usepackage{amsfonts}
\usepackage{amsmath}
\usepackage{amssymb}
\usepackage{graphicx}     

\begin{document}
  
\noindent \large Auteur: Abdellah El Moussaoui \\
\noindent \large Studierichting: Informatica\\
\noindent \large Oefeningenreeks 6B nr. 17.3\\

\medskip

\normalsize

$\displaystyle s_n=\sum_{k=1}^{n}k(k+1)(k+2)$\\

$= \displaystyle\sum_{k=1}^{n}k\left(k^2+3k+2\right)$\\

$= \displaystyle\sum_{k=1}^{n}k^3+3k^2+2k$\\

$= \displaystyle\sum_{k=1}^{n}k^3+3\displaystyle\sum_{k=1}^{n}k^2+2\displaystyle\sum_{k=1}^{n}k$\\

$= \dfrac{n^2\left(n+1\right)^2}{4}+3\cdot \dfrac{n\left(n+1\right)\left(2n+1\right)}{6}+ 2\cdot \dfrac{n\left(n+1\right)}{2}$\\

$= \dfrac{n^2\left(n+1\right)^2}{4}+\dfrac{n\left(n+1\right)\left(2n+1\right)}{2}+n\left(n+1\right)$\\

$= \dfrac{n\left(n+1\right)\left(n\left(n+1\right)+2\left(2n+1\right)+4\right)}{4}$\\

$= \dfrac{n\left(n+1\right)\left(n^2+n+4n+2+4\right)}{4}$\\

$= \dfrac{n\left(n+1\right)\left(n^2+5n+6\right)}{4}$\\


$= \dfrac{n\left(n+1\right)\left(n+2\right)\left(n+3\right)}{4}$\\


Dus de algemene term van $s_n$ is altijd geheel omdat: \\

$n\left(n+1\right)\left(n+2\right)\left(n+3\right) \mod 4 = 0$\\

\begin{thebibliography}{999}
%\bibitem{a} Referentie
\end{thebibliography}
\end{document} 